\documentclass[10pt,a4paper]{article}
\usepackage[margin=0.76in]{geometry}
\usepackage{amsmath,amssymb,booktabs,array}
\usepackage[hidelinks]{hyperref}
\setlength{\parindent}{0pt}
\setlength{\parskip}{4pt}
\newcommand{\head}[1]{\par\medskip\textbf{#1}\par}
\newcommand{\pos}[1]{\left[#1\right]_+}
\begin{document}
\begin{center}
{\large\bfseries Joined lane-drop VSL candidate: contract and qualification}\\[3pt]
2026-10-21 --- LF1/LF2 selected; concrete member delivered at JQ1
\end{center}
\textbf{Reading map.} Pages 1--2 specify the traffic member; pages 3--6 give
information, complete-accounting and derivative qualifications. These are
analytical results and declared defaults, not simulator or traffic experiments.
The earlier accepted VSL report is unchanged.

\head{1. Concrete member and global assumptions}
There are $N=10$ road rows of length $\Delta x=0.1$ km, each with lanes
$E,A,B$ in that left-to-right order. At the downstream boundary E ends;
A and B have exit outlets. Traffic completes when it crosses that line.
There are no physical downstream cells or recovery zone in this member.
The whole finite cohort is present at time zero; later external inflow is
zero. Its prescribed initial continuous stocks do not depend on the advice.
Lane stocks and incentives retain the principal framework~\cite{ns}, with
the explicit budget/allocation substitutions below.

All lanes use nominal free speed $u_0=60$ km/h, backward wave-speed magnitude
$w=20$ km/h, jam density $k_J=150$ veh/km and maximum capacity-drop fraction
$\delta=0.1$. Let $\Delta t=1/600$ h (6 s). One common constant advice
$v\in[30,60]$ km/h applies throughout the modeled road. These are illustrative
review defaults, not field calibration or an experimentally selected regime.
Use continuous vehicle-equivalent stock $n^k_{i\ell}=\Delta x\,\rho^k_{i\ell}$,
$0\leq n^0_{i\ell}\leq J=\Delta x k_J$; time index is $k$.

\head{2. Selected H budget and fixed boundaries}
Define the common FD quantities and Han recovery slope~\cite{han16,han17}:
\begin{align}
 C&=\frac{u_0wk_J}{u_0+w}=2250, & \rho_c&=C/u_0=37.5,\nonumber\\
 Q_J&=(1-\delta)C, & w_d&=\frac{Q_J}{k_J-Q_J/u_0}=\frac{540}{31}.
 \label{eq:fd}
\end{align}
In rates (veh/h), use
\begin{align}
 Q(x)&=\min\left\{C,C\left(1-\delta\frac{x-\rho_c}{k_J-\rho_c}\right)\right\},\nonumber\\
 R(x,y)&=\min\{Q(x),\ w(k_J-y),\ w(k_J-x)+w_d(x-y)\},\nonumber\\
 D_{i\ell}&=\min\{v\rho_{i\ell},Q(\rho_{i-1,\ell})\}. \label{eq:H}
\end{align}
Set $\rho_{0\ell}=0$ for the no-inflow upstream boundary. For a target lane
$j$ in row $i+1\leq N$, its \emph{one shared} receiving amount is
$r_{ij}=\Delta t R(\rho_{ij},\rho_{i+1,j})$. At $i=N$, only $j=A,B$ exist;
set $r_{Nj}=\Delta t R(\rho_{Nj},0)$ using free outlet ghost densities.
Ghosts are boundary values, not stocks or additional traffic. Advice alters
sending, not $Q,R$, priority weights, geometry or nominal $\rho_c$ directly.

All movements advance one row, diagonally or in the same lane. Evaluate
every flow from the old time level and then update all physical stocks
together. At interior boundaries allow every adjacent movement and the
same-lane movement. At $N$, E can only enter A; A/B may enter their own or
the other continuing outlet. This is the selected flow-realization adaptation,
not literal reuse of the principal paper's adjusted-flow formulas.

The selected H proof gives $0\leq D_{i\ell}\leq v\rho_{i\ell}$ and
$0\leq R(x,y)\leq w(k_J-y)$. Because $u_0\Delta t\leq\Delta x$ and
$w\Delta t\leq\Delta x$, demand amounts $d_{i\ell}=\Delta t D_{i\ell}$
cannot exceed stock and receiving amounts cannot exceed old available
storage. The allocator below therefore preserves nonnegative stocks,
storage, permissions and mass. Native calibration and wave fidelity do not
follow from these inequalities.

\newpage
\head{3. Explicit lane preferences and their local conventions}
Use the principal source's density/incentive framework~\cite{ns} and
same-author lookahead weights $(a_0,a_1,a_2)=(2,2,1)$~\cite{thesis}:
\begin{equation}
 K_{i\ell}=\frac{\sum_{r\in\mathcal H_{i\ell}}a_r\rho_{i+r,\ell}}
 {\sum_{r\in\mathcal H_{i\ell}}a_r},\quad
 P_{i,\ell j}=\pos{\frac{I_{i,\ell j}K_{i\ell}-K_{ij}}{K_{i\ell}+K_{ij}}},
 \quad I=1+I^r+I^{kr}+I^{coop}. \label{eq:pref}
\end{equation}
For A/B take $\mathcal H=\{0,1,2\}$ and use zero-density outlet values beyond
$N$. For E retain only $r$ with $i+r\leq N$ and renormalize the remaining
weights. This finite-domain average is a declared local convention; no
imaginary E congestion or extra mass is introduced. Empty donors have zero
requests without evaluating incentive ratios. A nonempty donor makes the
denominator in \eqref{eq:pref} positive.

Let $g_i=(N-i)\Delta x$, with $D_r=0.75$ km. Retain the displayed additive
source form, with route and courtesy terms only in the stated directions:
\begin{align}
 I^r_{i,EA}&=(1-g_i/D_r)^3\quad(0\leq g_i\leq D_r),\nonumber\\
 I^{coop}_{i,AB}&=(\rho_{iA}+\rho_{iB})/\rho_{iA}
 \quad(\rho_{iA}<\rho_c,\ 0\leq g_i<D_r). \label{eq:incentives}
\end{align}
These terms are zero otherwise. In free-flow origin cells
($0<\rho_{i\ell}<\rho_c$), use
$I^{kr}_{i,\ell j}=-\rho_{ij}/\rho_{i\ell}$ for a leftward move or a move
into rightmost B, and zero otherwise; disable it for the active E route move.
In congested origin cells use the source's allowed $I^{kr}=0$ choice. Equality
at $\rho_c$ belongs to that congested branch. The density-mode threshold is
nominal and independent of advice. These explicit direction/vacuum/tie and
domain conventions are part of the review member, not a recovered author
implementation. Route bias is not substituted for a realized merge fraction.

\head{4. Selected allocator, complete update and completion}
Mask nonexistent destinations before forming desires. For each nonterminal-E
donor, normalize its allowed lateral requests jointly:
\begin{equation}
 s_{i\ell}=\sum_{j\ne\ell}P_{i,\ell j},\qquad
 b_{i,\ell j}=\frac{P_{i,\ell j}}{\max(1,s_{i\ell})},\quad
 b_{i,\ell\ell}=1-\sum_{j\ne\ell}b_{i,\ell j},\quad
 m_{i,\ell j}=d_{i\ell}b_{i,\ell j}. \label{eq:request}
\end{equation}
At last-E set $b_{N,EA}=1$. All equal reference capacities give equal positive
priorities. For each target $j$, if requests fit its receiving amount, set
$z_{i,\ell j}=m_{i,\ell j}$. Otherwise choose the common level $\lambda_{ij}$
such that
\begin{equation}
 z_{i,\ell j}=\min(m_{i,\ell j},\lambda_{ij}),\qquad
 \sum_\ell z_{i,\ell j}=r_{ij}. \label{eq:allocate}
\end{equation}
This capped sharing returns unused entitlement~\cite{carey,wright}; its
three-lane composition is the selected adaptation. Give unused donor traffic
one forward fallback:
\begin{align}
 y_{i\ell}&=\min\left\{d_{i\ell}-\sum_jz_{i,\ell j},\ 
                r_{i\ell}-\sum_hz_{i,h\ell}\right\},\nonumber\\
 f_{i,\ell j}&=z_{i,\ell j}+\mathbf1_{\{j=\ell\}}y_{i\ell},\qquad
 n^{k+1}_{i\ell}=n^k_{i\ell}-\sum_jf^k_{i,\ell j}+\sum_hf^k_{i-1,h\ell}.
 \label{eq:update}
\end{align}
Set $y_{NE}=0$, $f_0=0$; terminal $f_N$ leaves the cohort and is recorded as
completion. No second lateral pass, extra headway multiplier, record history,
lossy compression or separate recovery meter is added. Each row sum of $f$
is at most its $d$, each target sum at most its $r$. The rule is deterministic
but need not maximize unrestricted total flow.

\newpage
\head{5. An actual same-observation lane-loading witness}
Use $v=60$. In row $N-1$, set all three lane densities to $\rho_c$; all
earlier rows are empty. Last-row densities are
\begin{equation}
 X:\ (\rho_E,\rho_A,\rho_B)=(120,60,90),\qquad
 Z:\ (\rho_E,\rho_A,\rho_B)=(120,90,60). \label{eq:states}
\end{equation}
The entire E profile and A+B profile coincide. Total initial mass is
$M_0=\Delta x(3\rho_c+120+60+90)=153/4$ veh-equivalents in either case.
The same observation equality holds for total row/link mass. These are
prescribed initial states; no identical past history is asserted.

All last-row donors have $Q(\rho_{N-1,\ell})=C$ and $D=C$. A/B are both
congested: courtesy is off and the declared congested keep-right term is
zero. Their route term is zero. The weighted A/B densities are proportional
to their current densities because both future outlet values are zero.
Thus the only continuing-lane lateral preference in X is $P_{BA}=1/5$;
in Z it is $P_{AB}=1/5$. E has the selected forced movement to A.

The receiving rates are
\begin{equation}
 (R_A,R_B)_X=(2205,2145),\qquad (R_A,R_B)_Z=(2145,2205).
\end{equation}
The selected allocation gives the following \emph{rates}; the update uses
these values multiplied by $\Delta t$:
\begin{center}
\begin{tabular}{lrrrr}
\toprule
 & $X:\to A$ & $X:\to B$ & $Z:\to A$ & $Z:\to B$\\
\midrule
E & 877.5 & 0 & 1072.5 & 0\\
A & 877.5 & 0 & 1072.5 & 450\\
B & 450 & 1800 & 0 & 1755\\
\midrule
Total & \multicolumn{2}{c}{4005} & \multicolumn{2}{c}{4350}\\
\bottomrule
\end{tabular}
\end{center}
For example in X, B's 450-unit request into A is fully served; E and A
share the remaining 1755 equally. B has used its whole donor budget, so
its unused outlet entitlement does not create more fallback. In Z both
receivers fill. E has no forward fallback in either case.

Their next total masses are therefore
\begin{equation}
 M_1(X)=\frac{153}{4}-\frac{4005}{600}=\frac{1263}{40},\qquad
 M_1(Z)=\frac{153}{4}-\frac{4350}{600}=31. \label{eq:nextmass}
\end{equation}
Hence the reduced observation does not close the actual service or next
aggregate evolution. The far-lane redistribution matters; an A-only monotonic
intuition is insufficient. This static one-step calculation is not a traffic
rollout or an optimal-control result.

For an equal-probability diagnostic pair, let $\mathcal J(v,x)$ be the full
objective below. Ideal reduced and full-information values are
\begin{align}
 V_{red}&=\inf_{v\in[30,60]}\tfrac12[\mathcal J(v,X)+\mathcal J(v,Z)],\nonumber\\
 V_{full}&=\tfrac12\inf_v\mathcal J(v,X)+\tfrac12\inf_v\mathcal J(v,Z),
 \qquad V_{red}-V_{full}\geq0. \label{eq:infovalue}
\end{align}
Strict positivity is unproved. One extra scalar $\rho_{NA}$ distinguishes X
and Z, so its augmented observer has exactly the full observer's ideal value
on this pair. Neither full lane detail nor useful different advice is thereby
necessary. Future aggregate flows reveal a distinction after the committed
decision; the claim is not absolute unobservability under histories/feedback.

\newpage
\head{6. Full population-time and eventual unconstrained service}
The objective is all remaining population-time from the initialized decision:
\begin{equation}
 M_k=\sum_{i,\ell}n^k_{i\ell},\qquad
 \mathcal J(v,n^0)=\Delta t\sum_{k=0}^{\infty}M_k. \label{eq:J}
\end{equation}
Its units are veh-equivalent hours. No later arrivals, uncounted holding
stock or residual deletion occurs. Trips complete at $f_N$, not upon merely
arriving at last-E. Action-independent pre-episode travel is outside this
remaining-time objective. Traffic after the free exit is outside the member.

Define remaining distance in cell units and a safe residual mass:
\begin{equation}
 h_i=N-i+1,\qquad V(n)=\sum_{i,\ell}h_i n_{i\ell},\qquad
 \eta=\Delta x\rho_c=15/4. \label{eq:potential}
\end{equation}
Each internal/exit movement advances one row and reduces $V$ by its mass.
If $\Gamma_k$ is the sum of all boundary movement amounts, including exits,
\begin{equation}
 V(n^{k+1})=V(n^k)-\Gamma_k,\qquad M_k\leq V(n^k)\leq N M_k.
 \label{eq:progress}
\end{equation}

\textbf{Why a nonempty state always makes progress.} Take its furthest
occupied row. Its receiving row is empty, or consists of exit outlets.
Every positive donor has positive $d$, because $v\geq30$ and $Q\geq Q_J>0$.
Every allowed receiver there has positive $r$ at zero receiving density.
Every donor has a permitted forward/adjacent destination, including E-to-A
at the endpoint. Positive equal priorities give a positive realized movement.
Thus $\Gamma>0$ whenever $V>0$, regardless of the desired fractions.

\textbf{Why this implies reaching small mass in finitely many steps.}
Allow the desired fractions $b$ to vary independently over the product of
their allowed simplices; this compact set contains every native incentive
output. For fixed $b$, the H budgets and capped-sharing/fallback flow amounts
are continuous in $(n,b,v)$, including zero requests. State bounds form a
compact set; advice belongs to $[30,60]$. On its closed subset $V\geq\eta$,
$\Gamma$ is positive everywhere, so it has a strictly positive minimum.
Equation~\eqref{eq:progress} therefore forces $V<\eta$ after finitely many
updates. This argument does not require the native incentive map itself to
be continuous, or an arbitrary convergence tolerance.

\textbf{Why every movement then receives service.} Once $M<\eta$, every
physical density and relevant upstream boundary density is below $\rho_c$.
Then $Q=C$ and $R=C$. Total requests into any target are at most
\begin{equation}
 \sum_\ell d_{i\ell}\leq\frac{u_0\Delta t}{\Delta x}M
 < C\Delta t=r_{ij}. \label{eq:tailspace}
\end{equation}
Thus $z=m$ at every target, no sending cap binds except $v\rho$, and each
physical donor advances exactly $c(v)n_{i\ell}$, where
\begin{equation}
 c(v)=v\Delta t/\Delta x\in[1/2,1].
\end{equation}
Mass cannot increase, so the regime persists. Consequently
$V^{k+1}=V^k-cM_k\leq(1-c/N)V^k$: the remaining mass decays geometrically.
Exact finite-time emptying is not required or claimed for all advice values.
The proof excludes continued arrivals, blocked exits, zero advice and a
changed transport/priority contract.

\newpage
\head{7. Exact residual value and full regular-branch derivative}
Choose any finite index $K$ inside the strict small-mass regime. Since
$V_k-V_{k+1}=cM_k$ there, telescope to get
\begin{equation}
 \Delta t\sum_{k=K}^{\infty}M_k
 =\frac{\Delta t}{c(v)}V(n^K)=\frac{\Delta x}{v}V(n^K).
 \label{eq:tail}
\end{equation}
This is the exact residual value of the discrete evaluator. A unit of mass
in row $i$ has mean remaining time $h_i\Delta x/v$; lane redistribution
cannot delay it when all requests receive service. It is not an approximate
terminal penalty or a rule deleting small stocks. Combining
\eqref{eq:J} and \eqref{eq:tail} gives the finite representation
\begin{equation}
 \mathcal J(v,n^0)=\Delta t\sum_{k=0}^{K-1}M_k+
                     \frac{\Delta x}{v}V(n^K). \label{eq:Jfinite}
\end{equation}
Later $K$ inside the same regime gives the same value. For example, with
$v=30$ a small unit of mass in last-E has mean remaining time
$\Delta t/c=0.1/30$ hours, including failed longitudinal advancement due to
the advice rate. A unit one row farther upstream takes twice that time.

Let the full state update in \eqref{eq:H}--\eqref{eq:update} be $F(n,v)$.
Suppose a neighborhood of advice preserves a regular finite branch trace
through time $K$. Let $F^{[k]}$ denote its selected smooth arithmetic formula
at step $k$, and write $s^k=\partial n^k/\partial v$. Then
\begin{equation}
 s^0=0,\qquad s^{k+1}=F^{[k]}_n(n^k,v)s^k+F^{[k]}_v(n^k,v), \label{eq:sens}
\end{equation}
and the derivative of \eqref{eq:Jfinite} is
\begin{equation}
 \frac{\partial\mathcal J}{\partial v}=
 \Delta t\sum_{k=0}^{K-1}\mathbf1^\top s^k+
 \frac{\Delta x}{v}\,h^\top s^K-
 \frac{\Delta x}{v^2}\,V(n^K). \label{eq:gradient}
\end{equation}
The three terms are finite-trajectory sensitivity, the residual stock's
sensitivity, and the direct change in residual travel time. The derivatives
$F^{[k]}_n,F^{[k]}_v$ include the source incentives, normalized requests, H budgets,
active receiving allocations, fallback and conservation update. They are
first derivatives of this complete new update, not the old report's service
function or a Hessian. Fixed initialization ties may use the selected branch
when all nearby advice follows that same choice; global differentiability
in arbitrary state directions is not assumed. Policy parameters can use
the subsequent chain rule through advice.

Choose $K$ with strict $M_K<\eta$ and hold that integer fixed for the local
derivative. A sufficiently small perturbation retains the exact terminal
representation. No derivative of a tolerance-based stopping index is
required. The tail already counts all future time and its dependence on $v$.
Differentiating only the finite trajectory would omit the last two terms.

\textbf{Regularity and events.} This result requires a stable finite sequence
of native incentive modes and active min/max/allocation constraints. Source
courtesy/keep-right changes can still be discontinuous at their state
thresholds. Empty donors and structural zero movements have their exact
mathematical treatment; no artificial density epsilon is assumed. At genuine
switches, two-sided derivatives may fail or finite action changes may cross
a jump. The evaluator uses a fixed time grid, so its regular derivative is
the discrete chain rule. Continuous-time saltation matrices from the old
point model do not apply to this different evaluator. No numerical gradient
agreement or useful policy direction has been tested.

\newpage
\head{8. Scientific limits, comparisons and provenance}
The observation ladder uses the same physical rows and forward law:
aggregate total; E plus A+B; and all lane states. Add the single adjacent
density statistic before attributing a need for all lane detail. The initial
pair in \eqref{eq:states} establishes service/transition nonclosure, not a
positive value in \eqref{eq:infovalue}. A model-aware reduced controller may
reason over the known hidden-state distribution. Histories or repeated
actions, if opened, require a fresh information comparison.

For the scalar action, matched-information bounded search and maximum advice
are mandatory strong references. A finite search does not certify a global
optimum. Separate information value, derivative correctness and decision
usefulness. Numerical derivative checks must include the complete residual
value; physical validation requires evidence beyond a decrease in this
same model's objective. SUMO is optional, as previously directed.

\textbf{Global assumptions supporting the new proofs.} Continuous mass,
fixed finite geometry and the common FD; closed initialized cohort and free
outlets; once-chosen positive common advice independent of initialization;
old-state conservative forward/diagonal movements with positive priorities;
completion at the exit and complete population-time; ordinary derivatives
only on regular traces. No microscopic order, driver records, integerization,
additional recovery law or claim of their necessity is introduced. Continuing
arrivals, heterogeneous FD mapping, blocked downstream traffic or different
control timing are not covered by these proofs automatically.

\textbf{Retained and new work.} Subraveti~\cite{ns} is the principal
lane-density/incentive framework. The same-author weights~\cite{thesis}
clarify its lookahead. Han~\cite{han16,han17}, already part of that congestion
lineage, gives the explicit budget/VSL component. Carey/Wright principles
\cite{carey,wright} support the selected conservative allocation adaptation.
The finite-domain conventions, common-FD member, information diagnostic and
complete-accounting/derivative arguments are explicitly our review work.
The selected composition is not a native published turnkey three-lane model.
No source calibration, continuum limit or maximal-throughput theorem transfers.

\textbf{Status.} LF1 and LF2 are user-selected. This bounded qualification is
complete for the declared member, delivered for confirmation at JQ1. It does
not replace the accepted report automatically, close old pending historical
gates, implement a simulator, or claim empirical/gradient success. No new
model-family choice is opened by this result.

\begin{thebibliography}{9}\setlength{\itemsep}{0pt}
\bibitem{ns} H. H. S. Nagalur Subraveti, V. L. Knoop, and B. van Arem,
``First order multi-lane traffic flow model -- an incentive based macroscopic
model to represent lane change dynamics,'' \emph{Transportmetrica B},
vol. 7, no. 1, pp. 1758--1779, doi: 10.1080/21680566.2019.1700846.
\bibitem{han16} Y. Han, Y. Yuan, A. Hegyi, and S. P. Hoogendoorn,
``New extended discrete first-order model to reproduce propagation of jam
waves,'' \emph{Transp. Res. Rec.}, vol. 2560, pp. 108--118, 2016,
doi: 10.3141/2560-12.
\bibitem{han17} Y. Han \emph{et al.}, ``Resolving freeway jam waves by
discrete first-order model-based predictive control of variable speed limits,''
\emph{Transp. Res. C}, vol. 77, pp. 405--420, 2017,
doi: 10.1016/j.trc.2017.02.009.
\bibitem{thesis} H. H. S. Nagalur Subraveti, \emph{Lane-specific Traffic
Flow Control}, Ph.D. dissertation, TU Delft, 2021,
doi: 10.4233/uuid:e405a549-b1fe-4fa1-8090-4429612359d9, chap. 4.
\bibitem{carey} M. Carey, C. Balijepalli, and D. Watling, ``Extending the cell
transmission model to multiple lanes and lane-changing,'' \emph{Netw. Spat.
Econ.}, vol. 15, pp. 507--535, 2015, doi: 10.1007/s11067-013-9193-7.
\bibitem{wright} M. A. Wright, G. Gomes, R. Horowitz, and A. A. Kurzhanskiy,
``On node models for high-dimensional road networks,'' \emph{Transp. Res. B},
vol. 105, pp. 212--234, 2017, doi: 10.1016/j.trb.2017.09.001.
\end{thebibliography}
\end{document}
